\section{A tensor obstruction detected by the scalar two-extensions}

\subsection{F-SM38: a countable scalar tensor cannot be exact in F-SM34}

Let \(\mathcal N\), \(E(V)=V^{**}/V\), and
\(\kappa(V)=(V,0)\) be as in F-SM34.  Suppose that
\(\mathcal N\) has an additive monoidal product with unit
\(\kappa(k)\), and that \(\kappa\) is strong monoidal:
\[
   \kappa(V)\otimes\kappa(W)\simeq\kappa(V\otimes_kW),
                                                               \tag{SM38.1}
\]
with the usual unit compatibility.  Then the functor
\[
       -\otimes\kappa(k^{(\mathbb N)}):
                   \mathcal N\longrightarrow\mathcal N       \tag{SM38.2}
\]
cannot be exact.  This uses only (SM38.1), not an
identification of the tensor of an arbitrary object with a
scalar object as a copower or as a componentwise tensor.
Closedness and symmetry are not needed.

The same conclusion holds for the external connective chain
category \(\mathcal M=\operatorname{Ch}_{\ge0}(\mathcal N)\)
and \(\kappa_{\mathcal M}(V)=\kappa(V)[0]\), for any monoidal
product with this unit and the analogous scalar compatibility.
Thus the actual additive interval example of F-SM34 cannot be
completed by a tensor which is exact even in this one countable
free scalar factor.  If the proposed tensor is right exact,
that scalar object is necessarily nonflat.

\begin{proof}
\emph{1. The canonical two-extension.}
Retain the functorial projective objects
\[
 P_0(V)=(V,V,1,0),\qquad P_1(V)=(E(V),V,0,1),
 \qquad S(V)=(0,V)
\]
from F-SM34.  There is a natural exact two-extension
\[
 \delta_V:\quad
 0\longrightarrow\kappa E(V)\longrightarrow P_1(V)
   \longrightarrow P_0(V)\longrightarrow\kappa(V)
   \longrightarrow0.                                      \tag{SM38.3}
\]
The middle map is zero in component zero and the identity
in component one.  Its kernel is \(\kappa E(V)\), its image
is \(S(V)\), and its cokernel is \(\kappa(V)\).

Dimension shifting through the two short exact sequences in
(SM38.3) gives natural identifications
\[
 \begin{aligned}
 \operatorname{Ext}^2_{\mathcal N}(\kappa V,\kappa U)
  &\simeq\operatorname{Ext}^1_{\mathcal N}(S(V),\kappa U)\\
  &\simeq\operatorname{Hom}_{\mathcal N}(\kappa E(V),\kappa U)
   \simeq\operatorname{Hom}_k(E(V),U).
 \end{aligned}                                               \tag{SM38.4}
\]
For the second identification,
\(\operatorname{Hom}_{\mathcal N}(P_1(V),\kappa U)=0\),
because \(P_1(V)\) represents the component-one functor,
whereas that component of \(\kappa U\) is zero.
The higher Ext groups of \(P_0(V)\) and \(P_1(V)\) vanish
by their projectivity.  These observations justify both
dimension-shifting isomorphisms, without any condition on
injective objects or infinite coproducts.  Under (SM38.4),
\(\delta_V\) is the identity of \(E(V)\): the connecting
map sends a homomorphism \(E(V)\to U\) to the pushout of
(SM38.3) along that homomorphism.  All identifications are
natural, contravariantly in \(V\) and covariantly in \(U\).

\emph{2. Exact tensoring would produce the prohibited contractions.}
Put \(W=k^{(\mathbb N)}\), and assume that (SM38.2) is
exact.  Tensoring (SM38.3) with \(\kappa W\) then produces
a genuine exact two-extension.  Using (SM38.1) and (SM38.4),
its class corresponds to a linear map
\[
          c:E(V\otimes W)\longrightarrow E(V)\otimes W.
                                                               \tag{SM38.5}
\]
For each coordinate functional \(\epsilon_n:W\to k\),
the natural map given by tensoring with \(\kappa(\epsilon_n)\)
is a morphism from this two-extension to \(\delta_V\).
The right unit isomorphism identifies the latter extension
with (SM38.3).  On the two endpoints these morphisms are
\(\kappa(1_{E(V)}\otimes\epsilon_n)\) and
\(\kappa(1_V\otimes\epsilon_n)\), respectively.
Naturality of (SM38.4), or equivalently equality of the
pushout and pullback classes of this morphism of extensions,
therefore gives
\[
    (1_{E(V)}\otimes\epsilon_n)c
                    =E(1_V\otimes\epsilon_n)\qquad(n\ge0).
                                                               \tag{SM38.6}
\]
Choose \(V\) and \(u\in E(V\otimes W)\) as in the complete
proof of F-SM36, so that the right side of (SM38.6), evaluated
at \(u\), is the same nonzero \(e\in E(V)\) for every
\(n\).  Then \(c(u)\) has infinitely many nonzero
coordinates in \(E(V)\otimes k^{(\mathbb N)}\), an
impossibility.  This proves the first assertion.  Only the
single functor (SM38.2) was assumed exact; tensoring with the
unit is automatically exact because it is isomorphic to the
identity.

\emph{3. Passage to the actual chain-category interval example.}
For a projective \(P\in\mathcal N\), the complex \(P[0]\)
is projective in \(\operatorname{Ch}_{\ge0}(\mathcal N)\):
maps from it are \(\operatorname{Hom}_{\mathcal N}(P,C_0)\),
and degree-zero evaluation is exact.  Applying \([0]\) to a
projective resolution in \(\mathcal N\) is a projective
resolution in this chain category, with exactness checked
degree by degree.  Consequently
\[
 \operatorname{Ext}^2_{\mathcal M}(\kappa V[0],\kappa U[0])
          \simeq\operatorname{Hom}_k(E(V),U),                 \tag{SM38.7}
\]
with the same universal class \(\delta_V[0]\).  The
argument of Step~2 applies verbatim to any tensor on
\(\mathcal M\) having the stipulated unit and scalar
compatibility.  No tensor inherited from \(\mathcal N\)
was assumed in this last assertion.
\end{proof}

\paragraph{Relation to the original question.}
The handwritten assumptions do not explicitly require
exactness of (SM38.2).  Thus F-SM38 is a stronger obstruction
to constructing a tensor on F-SM34, and it does not decide
the remaining abstract bilinear-tensor question.  Its input
is scalar strong-monoidal compatibility alone; it is distinct
from F-SM35's argument using exact copowers and products.
